\begin{abstract}
Support Vector Clustering (SVC) can theoretically capture cluster boundaries of arbitrary shape by seeking a minimal enclosing hypersphere in a high-dimensional feature space. However, its clustering performance exhibits extreme ``cliff-like'' sensitivity to the Gaussian kernel width parameter $q$: a slight shift in $q$ can cause the model to collapse from perfect clustering (ARI~=~1.000) to complete failure (ARI~=~0.000). This parameter fragility severely limits the deployment of SVC in real-world unsupervised scenarios---without ground-truth labels, neither cross-validation nor grid search can determine the optimal $q$. To address this problem, we propose Adaptive Support Vector Clustering (ASVC), which makes three core contributions: (1)~an automatic $q$-range estimation algorithm based on $k$-NN local neighborhood distances, which uses local geometric statistics in place of global pairwise distances to fundamentally avoid cross-cluster contamination; (2)~a two-phase ``coarse-to-fine'' parameter search strategy, which first explores on a logarithmic scale and then refines linearly around the coarse optimum; and (3)~a dedicated continuous scoring function combining a cluster-count bonus, a support vector ratio preference, and a piecewise bounded support vector (BSV) penalty, making candidate solutions across the parameter space distinguishable and comparable. Systematic experiments on six synthetic datasets covering diverse geometric structures show that ASVC maintains a baseline ARI~$\geq$~0.428 across all datasets, and reverses SVC's complete failure (ARI~=~0) to ARI~=~0.960--0.985 on two difficult datasets where the original SVC fails entirely. Notably, on our constructed VaryingBlobs dataset, ASVC demonstrates density-independence, outperforming the density-based DBSCAN. Our core claim is that ASVC transforms SVC from a ``laboratory method'' requiring ground-truth guidance for parameter tuning into a practical clustering tool that can operate independently in unsupervised settings.

All experimental code is open-sourced at \url{https://github.com/tuantuan-zi/ASVC}.
\end{abstract}
